% ============================================================
%  CHAPITRE 4 : IMPLÉMENTATION
% ============================================================
\chapter{Implémentation}

\section*{Introduction}
Ce chapitre décrit l'implémentation concrète du système d'anonymisation d'images médicales. Il présente l'environnement de développement, la structure du projet, puis détaille les extraits de code les plus significatifs de chaque module~: backend Node.js, pipeline IA FastAPI, stockage MinIO et interface React.

\section{Environnement de développement}

\subsection{Outils et technologies}

L'environnement de développement repose sur les outils suivants~: Visual Studio Code comme éditeur principal, Git et GitHub pour la gestion de versions et la collaboration, et un environnement virtuel Python \texttt{venv} pour l'isolation des dépendances. Le développement a été réalisé sur Windows~11 avec Node.js~v22.20.0, Python~3.10, MongoDB~8.0 local, et MinIO déployé via Docker.

\subsection{Structure complète du projet}

\begin{lstlisting}[language=bash, caption={Arborescence du projet PFE\_Test}, label={lst:structure}]
PFE_Test/
|-- api/                        # Pipeline IA FastAPI
|   |-- main.py                 # Point d'entree + 7 etapes du pipeline
|   |-- config.py               # Configuration MinIO (Pydantic Settings)
|   +-- storage.py              # Client MinIO S3-compatible
|-- anonymizer/                 # Modules d'anonymisation
|   |-- pixel_redactor.py       # Redaction adaptative OpenCV TELEA
|   |-- metadata_anonymizer.py  # Suppression des tags DICOM PHI
|   |-- image_validator.py      # Validation format + extraction pixels
|   +-- dicom_decompressor.py   # Decompression DICOM (gdcm/pylibjpeg)
|-- ocr/                        # Moteurs OCR
|   |-- text_detector.py        # PaddleOCR PP-OCRv4
|   |-- easy_text_detector.py   # EasyOCR (region bordur.)
|   +-- preprocessor.py         # CLAHE - BorderPreprocessor
|-- image_classifier/
|   +-- improved_medical_classifier.py  # CLIP ViT-B/32 zero-shot
|-- backend/                    # Node.js Express API
|   |-- src/
|   |   |-- app.js              # Configuration Express + CORS + routes
|   |   |-- server.js           # Point d'entree + connexion MongoDB
|   |   |-- controllers/
|   |   |   |-- authController.js       # Inscription, connexion, profil
|   |   |   |-- anonymizeController.js  # Orchestration pipeline IA
|   |   |   +-- historyController.js    # Historique + admin
|   |   |-- middleware/
|   |   |   |-- auth.js         # JWT protect + RBAC middlewares
|   |   |   +-- upload.js       # Multer memoryStorage
|   |   |-- models/
|   |   |   |-- User.js         # Schema Mongoose utilisateur
|   |   |   +-- History.js      # Schema Mongoose journal
|   |   +-- routes/
|   |       |-- auth.js
|   |       |-- anonymize.js
|   |       +-- history.js
|   |-- package.json
|   +-- .env                    # Variables d'environnement
|-- client/                     # React + Vite
|   +-- src/
|       |-- App.jsx             # Router + AuthContext provider
|       |-- pages/
|       |   |-- Login.jsx       # Page connexion/inscription
|       |   |-- Dashboard.jsx   # Tableau de bord principal
|       |   +-- AdminPanel.jsx  # Panneau administrateur
|       |-- components/
|       |   |-- UploadZone.jsx  # Zone upload + paramètres avancés
|       |   |-- AdvancedSettings.jsx  # Curseurs paramètres pipeline
|       |   +-- HistoryList.jsx # Historique des operations
|       +-- context/
|           +-- AuthContext.jsx # Contexte auth global
|-- requirements.txt            # Dependances Python fixees
|-- docker-compose.yml          # MinIO + services conteneurises
+-- .gitignore
\end{lstlisting}

\section{Implémentation du backend Node.js}

\subsection{Configuration de l'application Express}

\begin{lstlisting}[language=JavaScript, caption={Configuration principale de l'application Express (app.js)}, label={lst:express}]
// backend/src/app.js
const express = require('express')
const cors = require('cors')
const morgan = require('morgan')
const app = express()

app.use(cors({
  origin: process.env.FRONTEND_URL || 'http://localhost:3000',
  credentials: true
}))
app.use(express.json({ limit: '10mb' }))
app.use(express.urlencoded({ extended: true }))
app.use(morgan('combined'))

// Routes
app.use('/api/auth',      require('./routes/auth'))
app.use('/api/anonymize', require('./routes/anonymize'))
app.use('/api/history',   require('./routes/history'))
app.use('/api/health', (req, res) =>
  res.json({ status: 'ok', timestamp: new Date() }))

// Gestion globale des erreurs
app.use((err, req, res, next) => {
  console.error(err.stack)
  res.status(500).json({ success: false, message: 'Internal Server Error' })
})

module.exports = app
\end{lstlisting}

\subsection{Module d'authentification}

\subsubsection{Modèle Mongoose User}

\begin{lstlisting}[language=JavaScript, caption={Modèle Mongoose User avec hachage bcrypt}, label={lst:user_model}]
// backend/src/models/User.js
const mongoose = require('mongoose')
const bcrypt   = require('bcryptjs')

const UserSchema = new mongoose.Schema({
  name: {
    type: String, required: [true, 'Nom requis'], maxlength: 50
  },
  email: {
    type: String, required: true,
    unique: true,
    match: [/^\S+@\S+\.\S+$/, 'Email invalide']
  },
  password: {
    type: String, required: true,
    minlength: [6, 'Minimum 6 caractères'],
    select: false   // jamais renvoyé dans les requêtes
  },
  role: {
    type: String,
    enum: ['utilisateur', 'utilisateur_medical', 'responsable'],
    default: 'utilisateur'
  }
}, { timestamps: true })

// Hachage automatique avant sauvegarde
UserSchema.pre('save', async function(next) {
  if (!this.isModified('password')) return next()
  this.password = await bcrypt.hash(this.password, 12)
  next()
})

UserSchema.methods.matchPassword = async function(entered) {
  return await bcrypt.compare(entered, this.password)
}

module.exports = mongoose.model('User', UserSchema)
\end{lstlisting}

\subsubsection{Contrôleur d'authentification}

\begin{lstlisting}[language=JavaScript, caption={Contrôleur d'authentification -- inscription et connexion}, label={lst:auth_ctrl}]
// backend/src/controllers/authController.js
const jwt  = require('jsonwebtoken')
const User = require('../models/User')

const generateToken = (id) =>
  jwt.sign({ id }, process.env.JWT_SECRET,
           { expiresIn: process.env.JWT_EXPIRE || '7d' })

// POST /api/auth/register
exports.register = async (req, res) => {
  try {
    const { name, email, password, role, adminKey } = req.body

    if (role === 'responsable') {
      if (adminKey !== process.env.ADMIN_REGISTRATION_KEY)
        return res.status(403).json({
          success: false, message: 'Clé admin invalide'
        })
    }
    const user = await User.create({ name, email, password, role })
    res.status(201).json({
      success: true,
      token: generateToken(user._id),
      user: { id: user._id, name: user.name,
              email: user.email, role: user.role }
    })
  } catch (err) {
    res.status(400).json({ success: false, message: err.message })
  }
}

// POST /api/auth/login
exports.login = async (req, res) => {
  try {
    const { email, password } = req.body
    const user = await User.findOne({ email }).select('+password')
    if (!user || !(await user.matchPassword(password)))
      return res.status(401).json({
        success: false, message: 'Identifiants invalides'
      })
    res.json({
      success: true,
      token: generateToken(user._id),
      user: { id: user._id, name: user.name,
              email: user.email, role: user.role }
    })
  } catch (err) {
    res.status(500).json({ success: false, message: err.message })
  }
}
\end{lstlisting}

\subsubsection{Middleware JWT et RBAC}

\begin{lstlisting}[language=JavaScript, caption={Middleware d'authentification JWT et contrôle RBAC}, label={lst:middleware_auth}]
// backend/src/middleware/auth.js
const jwt  = require('jsonwebtoken')
const User = require('../models/User')

exports.protect = async (req, res, next) => {
  try {
    let token
    if (req.headers.authorization?.startsWith('Bearer'))
      token = req.headers.authorization.split(' ')[1]
    if (!token)
      return res.status(401).json({ success: false,
        message: 'Non autorisé -- token manquant' })

    const decoded = jwt.verify(token, process.env.JWT_SECRET)
    req.user = await User.findById(decoded.id)
    if (!req.user)
      return res.status(401).json({ success: false,
        message: 'Utilisateur introuvable' })
    next()
  } catch {
    res.status(401).json({ success: false,
      message: 'Non autorisé -- token invalide' })
  }
}

exports.responsableOnly = (req, res, next) => {
  if (req.user?.role !== 'responsable')
    return res.status(403).json({ success: false,
      message: 'Accès réservé aux responsables' })
  next()
}

exports.medicalOrResponsable = (req, res, next) => {
  const allowed = ['utilisateur_medical', 'responsable']
  if (!allowed.includes(req.user?.role))
    return res.status(403).json({ success: false,
      message: 'Accès réservé aux utilisateurs médicaux' })
  next()
}
\end{lstlisting}

\subsection{Contrôleur d'anonymisation}

\begin{lstlisting}[language=JavaScript, caption={Contrôleur d'anonymisation -- transmission au pipeline IA}, label={lst:anon_ctrl}]
// backend/src/controllers/anonymizeController.js
const axios    = require('axios')
const FormData = require('form-data')
const { Blob } = require('buffer')
const History  = require('../models/History')

exports.anonymizeImage = async (req, res) => {
  const startTime = Date.now()
  try {
    // Extraction des paramètres du pipeline
    const conf_threshold = parseFloat(req.body.conf_threshold) || 0.1
    const padding        = parseInt(req.body.padding)          || 5
    const border_margin  = parseInt(req.body.border_margin)    || 100
    const border_pct     = parseFloat(req.body.border_pct)     || 0.20

    // Construction du FormData pour FastAPI
    const formData = new FormData()
    const blob = new Blob([req.file.buffer],
                          { type: req.file.mimetype })
    formData.append('file',          blob, req.file.originalname)
    formData.append('conf_threshold',conf_threshold.toString())
    formData.append('padding',       padding.toString())
    formData.append('border_margin', border_margin.toString())
    formData.append('border_pct',    border_pct.toString())

    // Appel au pipeline FastAPI
    const response = await axios.post(
      `${process.env.FASTAPI_URL}/anonymize`,
      formData,
      { headers: formData.getHeaders(), timeout: 120000 }
    )
    const data = response.data

    // Persistance du journal
    await History.create({
      user: req.user._id,
      filename: req.file.originalname,
      status: 'success',
      category: data.category,
      confidence: data.confidence,
      regionsDetected: data.regions_detected,
      regionsRedacted: data.regions_redacted,
      tagsAnonymized: data.tags_anonymized,
      processingTime: Date.now() - startTime,
      minioUri: data.minio_uri,
      downloadUrl: data.download_url,
      confThreshold: conf_threshold,
      padding, borderMargin: border_margin, borderPct: border_pct
    })

    res.json({ success: true, data })
  } catch (err) {
    await History.create({
      user: req.user._id,
      filename: req.file?.originalname || 'unknown',
      status: 'failed',
      processingTime: Date.now() - startTime
    })
    res.status(500).json({ success: false, message: err.message })
  }
}
\end{lstlisting}

\section{Implémentation du pipeline IA (FastAPI)}

\subsection{Point d'entrée FastAPI et configuration CORS}

\begin{lstlisting}[language=Python, caption={Configuration FastAPI et CORS (main.py -- extrait)}, label={lst:fastapi_config}]
# api/main.py (extrait)
from fastapi import FastAPI, File, UploadFile, Form, HTTPException
from fastapi.middleware.cors import CORSMiddleware
from fastapi.responses import JSONResponse
import uvicorn, logging, numpy as np
from datetime import datetime
from pathlib import Path
from PIL import Image

app = FastAPI(
    title="Medical Image Anonymization API",
    description="7-stage AI anonymization pipeline",
    version="2.0.0"
)

app.add_middleware(CORSMiddleware,
    allow_origins=["http://localhost:3000",
                   "http://localhost:3001",
                   "http://localhost:5000"],
    allow_credentials=True,
    allow_methods=["*"],
    allow_headers=["*"]
)

TEMP_DIR   = Path("temp")
OUTPUT_DIR = Path("output")
TEMP_DIR.mkdir(exist_ok=True)
OUTPUT_DIR.mkdir(exist_ok=True)
logger = logging.getLogger(__name__)
\end{lstlisting}

\subsection{Classifieur CLIP zéro-shot}

\begin{lstlisting}[language=Python, caption={Classifieur CLIP ViT-B/32 zéro-shot}, label={lst:clip}]
# image_classifier/improved_medical_classifier.py (extrait)
from transformers import CLIPProcessor, CLIPModel
from PIL import Image
import torch

class MedicalImageClassifier:
    CATEGORIES = {
        "chest":        "a chest X-ray radiograph image",
        "dental":       "a dental X-ray image",
        "pelvic":       "a pelvic X-ray scan image",
        "skull":        "a skull X-ray radiograph image",
        "other_medical":"a medical scan other than chest or dental",
        "non_medical":  "a non-medical photograph or document"
    }

    def __init__(self):
        self.model = CLIPModel.from_pretrained("openai/clip-vit-base-patch32")
        self.processor = CLIPProcessor.from_pretrained(
                             "openai/clip-vit-base-patch32")

    def classify_image(self, image_path: str):
        image = Image.open(image_path).convert("RGB")
        texts = list(self.CATEGORIES.values())
        inputs = self.processor(text=texts, images=image,
                                return_tensors="pt", padding=True)
        with torch.no_grad():
            outputs = self.model(**inputs)
        probs = outputs.logits_per_image.softmax(dim=1)[0]
        best_idx = probs.argmax().item()
        category = list(self.CATEGORIES.keys())[best_idx]
        confidence = float(probs[best_idx])
        return category, confidence, {}
\end{lstlisting}

\subsection{Préprocesseur CLAHE bordural}

\begin{lstlisting}[language=Python, caption={BorderPreprocessor -- CLAHE sur les régions bordurales}, label={lst:clahe}]
# ocr/preprocessor.py (extrait)
import cv2, numpy as np

class BorderPreprocessor:
    def __init__(self, border_pct: float = 0.15):
        self.border_pct = border_pct
        self.clahe = cv2.createCLAHE(
            clipLimit=2.0, tileGridSize=(8, 8))

    def enhance(self, image: np.ndarray) -> np.ndarray:
        """Applique CLAHE uniquement sur les bandes bordurales."""
        H, W = image.shape[:2]
        margin_h = max(1, int(H * self.border_pct))
        margin_w = max(1, int(W * self.border_pct))
        result = image.copy()

        def apply_clahe(region):
            if len(region.shape) == 3:
                lab = cv2.cvtColor(region, cv2.COLOR_RGB2LAB)
                lab[:,:,0] = self.clahe.apply(lab[:,:,0])
                return cv2.cvtColor(lab, cv2.COLOR_LAB2RGB)
            return self.clahe.apply(region)

        # Appliquer sur les 4 bandes
        for slc in [
            (slice(0, margin_h), slice(None)),          # haut
            (slice(H-margin_h, H), slice(None)),         # bas
            (slice(None), slice(0, margin_w)),           # gauche
            (slice(None), slice(W-margin_w, W))          # droite
        ]:
            result[slc] = apply_clahe(image[slc])
        return result
\end{lstlisting}

\subsection{Module de redaction adaptatif}

\begin{lstlisting}[language=Python, caption={PixelRedactor -- redaction adaptative avec échantillonnage de fond}, label={lst:redactor}]
# anonymizer/pixel_redactor.py (extrait)
import cv2, numpy as np

class PixelRedactor:
    def __init__(self, padding: int = 5, border_margin: int = 100):
        self.padding = padding
        self.border_margin = border_margin

    def _is_on_black_frame(self, pixels, x_min, y_min, x_max, y_max,
                           black_threshold: int = 60) -> bool:
        H, W = pixels.shape[:2]
        s = pixels[max(0,y_min-15):min(H,y_max+15),
                   max(0,x_min-15):min(W,x_max+15)]
        if len(s.shape) == 3:
            s = s.mean(axis=2)
        return float(np.mean(s)) < black_threshold

    def _sample_edge_fill(self, pixels, x_min, y_min, x_max, y_max):
        """Echantillonne la couleur de fond depuis le bord le plus proche."""
        H, W = pixels.shape[:2]
        if x_min < W - x_max:  # plus proche du bord gauche
            strip = pixels[y_min:y_max, 0:min(15, W)]
        else:                   # plus proche du bord droit
            strip = pixels[y_min:y_max, max(0, W-15):W]
        if len(pixels.shape) == 3:
            return [int(np.median(strip[:,:,c])) for c in range(3)]
        return int(np.median(strip))

    def _redact_region(self, pixels, x_min, y_min, x_max, y_max,
                       padding, is_grayscale):
        H, W = pixels.shape[:2]
        xp0 = max(0, x_min - padding)
        yp0 = max(0, y_min - padding)
        xp1 = min(W, x_max + padding)
        yp1 = min(H, y_max + padding)

        if self._is_on_black_frame(pixels, x_min, y_min, x_max, y_max):
            fill = self._sample_edge_fill(pixels, x_min, y_min, x_max, y_max)
            pixels[yp0:yp1, xp0:xp1] = fill
        else:
            # Inpainting TELEA rayon 1
            mask = np.zeros(pixels.shape[:2], dtype=np.uint8)
            mask[yp0:yp1, xp0:xp1] = 255
            if is_grayscale:
                pixels = cv2.inpaint(pixels, mask,
                                     inpaintRadius=1,
                                     flags=cv2.INPAINT_TELEA)
            else:
                channels = []
                for c in range(pixels.shape[2]):
                    ch = cv2.inpaint(pixels[:,:,c], mask,
                                     inpaintRadius=1,
                                     flags=cv2.INPAINT_TELEA)
                    channels.append(ch)
                pixels = np.stack(channels, axis=2)
        return pixels
\end{lstlisting}

\subsection{Détection OCR double moteur avec fusion IoU}

\begin{lstlisting}[language=Python, caption={Fusion des résultats PaddleOCR et EasyOCR par IoU}, label={lst:iou_fusion}]
# api/main.py -- fonction merge_and_deduplicate (extrait)
def merge_and_deduplicate(paddle_regions, easy_regions,
                          iou_threshold=0.5):
    def compute_iou(r1, r2):
        x1 = [p[0] for p in r1]; y1 = [p[1] for p in r1]
        x2 = [p[0] for p in r2]; y2 = [p[1] for p in r2]
        inter_x = max(0, min(max(x1),max(x2))-max(min(x1),min(x2)))
        inter_y = max(0, min(max(y1),max(y2))-max(min(y1),min(y2)))
        if inter_x == 0 or inter_y == 0:
            return 0.0
        inter  = inter_x * inter_y
        area1  = (max(x1)-min(x1)) * (max(y1)-min(y1))
        area2  = (max(x2)-min(x2)) * (max(y2)-min(y2))
        return inter / (area1 + area2 - inter)

    merged = list(paddle_regions)
    for easy in easy_regions:
        if not any(compute_iou(easy, p) > iou_threshold
                   for p in paddle_regions):
            merged.append(easy)
    return merged
\end{lstlisting}

\section{Implémentation du stockage MinIO}

\begin{lstlisting}[language=Python, caption={Client MinIO -- upload et URL pré-signée}, label={lst:minio}]
# api/storage.py (extrait)
from minio import Minio
from minio.error import S3Error
from datetime import timedelta
import os

class MinIOStorage:
    def __init__(self, endpoint, access_key, secret_key,
                 bucket_name, secure=False):
        self.client = Minio(endpoint, access_key=access_key,
                            secret_key=secret_key, secure=secure)
        self.bucket_name = bucket_name
        self._ensure_bucket()

    def _ensure_bucket(self):
        if not self.client.bucket_exists(self.bucket_name):
            self.client.make_bucket(self.bucket_name)

    def upload_file(self, file_path: str, category: str = "other"):
        from datetime import datetime
        now = datetime.now()
        filename = os.path.basename(file_path)
        object_name = (f"{category}/{now.year}/{now.month:02d}/"
                       f"{now.day:02d}/{filename}")
        self.client.fput_object(
            self.bucket_name, object_name, file_path)
        return f"minio://{self.client._base_url.netloc}/{self.bucket_name}/{object_name}"

    def get_presigned_url(self, object_name: str,
                          expires_hours: int = 24) -> str:
        return self.client.presigned_get_object(
            self.bucket_name, object_name,
            expires=timedelta(hours=expires_hours))
\end{lstlisting}

\section{Implémentation de l'interface React}

\subsection{Contexte d'authentification global}

\begin{lstlisting}[language=JavaScript, caption={AuthContext React -- gestion de l'état d'authentification}, label={lst:authcontext}]
// client/src/context/AuthContext.jsx (extrait)
import { createContext, useContext, useState } from 'react'
import axios from 'axios'

const AuthContext = createContext()
const API = 'http://localhost:5000/api'

export const AuthProvider = ({ children }) => {
  const [user, setUser] = useState(
    JSON.parse(localStorage.getItem('user') || 'null'))
  const [token, setToken] = useState(
    localStorage.getItem('token') || null)

  const login = async (email, password) => {
    const { data } = await axios.post(`${API}/auth/login`,
                                      { email, password })
    setUser(data.user); setToken(data.token)
    localStorage.setItem('user', JSON.stringify(data.user))
    localStorage.setItem('token', data.token)
  }

  const logout = () => {
    setUser(null); setToken(null)
    localStorage.clear()
  }

  // Intercepteur axios pour ajouter le token JWT
  axios.defaults.headers.common['Authorization'] =
    token ? `Bearer ${token}` : ''

  return (
    <AuthContext.Provider value={{ user, token, login, logout }}>
      {children}
    </AuthContext.Provider>
  )
}

export const useAuth = () => useContext(AuthContext)
\end{lstlisting}

\subsection{Composant AdvancedSettings -- Paramètres du pipeline}

\begin{lstlisting}[language=JavaScript, caption={Composant AdvancedSettings -- curseurs de paramétrage du pipeline}, label={lst:advsettings}]
// client/src/components/AdvancedSettings.jsx (extrait)
import React from 'react'

export default function AdvancedSettings({ settings, onChange }) {
  const params = [
    { key: 'conf_threshold', label: 'Seuil de confiance OCR',
      min: 0.0, max: 1.0, step: 0.05,
      hint: 'Valeur basse = plus de détections' },
    { key: 'padding',        label: 'Rembourrage (px)',
      min: 0, max: 20, step: 1 },
    { key: 'border_margin',  label: 'Marge de sécurité (px)',
      min: 50, max: 300, step: 10 },
    { key: 'border_pct',     label: 'Zone bordur. EasyOCR (%)',
      min: 0.05, max: 0.50, step: 0.05 }
  ]
  return (
    <div className="advanced-settings">
      {params.map(({ key, label, min, max, step, hint }) => (
        <div key={key} className="param-row">
          <label>{label}: <strong>{settings[key]}</strong></label>
          <input type="range" min={min} max={max} step={step}
            value={settings[key]}
            onChange={e => onChange(key,
              step < 1 ? parseFloat(e.target.value)
                       : parseInt(e.target.value))}
          />
          {hint && <span className="hint">{hint}</span>}
        </div>
      ))}
    </div>
  )
}
\end{lstlisting}

\section{Configuration DevOps}

\subsection{Variables d'environnement backend}

\begin{lstlisting}[language=bash, caption={Fichier .env du backend (valeurs d'exemple)}, label={lst:env}]
# backend/.env
PORT=5000
MONGO_URI=mongodb://localhost:27017/medical_anonymizer
JWT_SECRET=votre_cle_secrete_jwt_tres_longue_et_aleatoire
JWT_EXPIRE=7d
FASTAPI_URL=http://localhost:8000
ADMIN_REGISTRATION_KEY=PURA_ADMIN_2026
FRONTEND_URL=http://localhost:3000
MINIO_ENDPOINT=localhost:9000
MINIO_ACCESS_KEY=minioadmin
MINIO_SECRET_KEY=minioadmin
MINIO_BUCKET=anonymized-images
\end{lstlisting}

\subsection{Docker Compose pour MinIO}

\begin{lstlisting}[language=bash, caption={docker-compose.yml -- service MinIO}, label={lst:docker}]
# docker-compose.yml
version: '3.8'
services:
  minio:
    image: minio/minio:latest
    container_name: minio
    ports:
      - "9000:9000"
      - "9001:9001"
    environment:
      MINIO_ROOT_USER: minioadmin
      MINIO_ROOT_PASSWORD: minioadmin
    volumes:
      - minio_data:/data
    command: server /data --console-address ":9001"
    restart: unless-stopped
volumes:
  minio_data:
\end{lstlisting}

\section*{Conclusion}
Ce chapitre a présenté l'implémentation concrète de l'ensemble des modules du système. Les extraits de code illustrent les choix d'implémentation pour chaque composant~: authentification JWT, pipeline IA en sept étapes, redaction adaptative, stockage MinIO et interface React. Le chapitre suivant présente la stratégie de tests et les résultats de validation obtenus.
